%% 
%% Copyright 2019-2024 Elsevier Ltd
%% 
%% This file is part of the 'CAS Bundle'.
%% --------------------------------------
%% 
%% It may be distributed under the conditions of the LaTeX Project Public
%% License, either version 1.3c of this license or (at your option) any
%% later version.  The latest version of this license is in
%%    http://www.latex-project.org/lppl.txt
%% and version 1.3c or later is part of all distributions of LaTeX
%% version 1999/12/01 or later.
%% 
%% The list of all files belonging to the 'CAS Bundle' is
%% given in the file `manifest.txt'.
%% 
%% Template article for cas-dc documentclass for 
%% double column output.

\documentclass[a4paper,fleqn]{cas-dc}

% If the frontmatter runs over more than one page
% use the longmktitle option.

%\documentclass[a4paper,fleqn,longmktitle]{cas-dc}

%\usepackage[numbers]{natbib}
%\usepackage[authoryear]{natbib}
\usepackage[authoryear,longnamesfirst]{natbib}

\usepackage{algorithm}
\usepackage{algorithmic}



%%%Author macros
\def\tsc#1{\csdef{#1}{\textsc{\lowercase{#1}}\xspace}}
\tsc{WGM}
\tsc{QE}
%%%

% Uncomment and use as if needed
%\newtheorem{theorem}{Theorem}
%\newtheorem{lemma}[theorem]{Lemma}
%\newdefinition{rmk}{Remark}
%\newproof{pf}{Proof}
%\newproof{pot}{Proof of Theorem \ref{thm}}

\begin{document}
\let\WriteBookmarks\relax
\def\floatpagepagefraction{1}
\def\textpagefraction{.001}

% Short title
\shorttitle{}    

% Short author
\shortauthors{}  

% Main title of the paper
\title [mode = title]{}  

% Title footnote mark
% eg: \tnotemark[1]
\tnotemark[1] 

% Title footnote 1.
% eg: \tnotetext[1]{Title footnote text}
\tnotetext[1]{} 

% First author
%
% Options: Use if required
% eg: \author[1,3]{Author Name}[type=editor,
%       style=chinese,
%       auid=000,
%       bioid=1,
%       prefix=Sir,
%       orcid=0000-0000-0000-0000,
%       facebook=<facebook id>,
%       twitter=<twitter id>,
%       linkedin=<linkedin id>,
%       gplus=<gplus id>]

\author[1]{}%[<options>]

% Corresponding author indication
\cormark[1]

% Footnote of the first author
\fnmark[1]

% Email id of the first author
\ead{}

% URL of the first author
\ead[url]{}

% Credit authorship
% eg: \credit{Conceptualization of this study, Methodology, Software}
\credit{}

% Address/affiliation
\affiliation[1]{organization={},
            addressline={}, 
            city={},
%          citysep={}, % Uncomment if no comma needed between city and postcode
            postcode={}, 
            state={},
            country={}}

\author[2]{}%[]

% Footnote of the second author
\fnmark[2]

% Email id of the second author
\ead{}

% URL of the second author
\ead[url]{}

% Credit authorship
\credit{}

% Address/affiliation
\affiliation[2]{organization={},
            addressline={}, 
            city={},
%          citysep={}, % Uncomment if no comma needed between city and postcode
            postcode={}, 
            state={},
            country={}}

% Corresponding author text
\cortext[1]{Corresponding author}

% Footnote text
\fntext[1]{}

% For a title note without a number/mark
%\nonumnote{}

% Here goes the abstract
\begin{abstract}
Here goes the abstract \nocite{*}%% Remove this line from your manuscript.
\end{abstract}

% Use if graphical abstract is present
%\begin{graphicalabstract}
%\includegraphics{}
%\end{graphicalabstract}

% Research highlights
\begin{highlights}
\item 
\item 
\item 
\end{highlights}


%\nocite{*}

% Keywords
% Each keyword is seperated by \sep
\begin{keywords}
 \sep \sep \sep
\end{keywords}

\maketitle

% Main text
\section{Introduction}

% Sta ovde mora da stoji:
% - Ne pisati kao digital twin paper.
% - Ne pisati kao plant biology experiment.
% - Prvi pasus: modelling/simulation problem.
% - Drugi pasus: zasto je relevantno za SMP&T.
% - Treci pasus: sta fali u postojecoj literaturi.
% - Cetvrti pasus: sta je novo u ovom radu.
% - Na kraju eksplicitni contributions.

\subsection{Motivation and modelling problem}
% Treba reci:
% - Simulacija rasta zahteva spajanje environment inputa,
%   gas exchange, carbon allocation, nutrients, biomass.
% - Photoperiod je dobar controlled scenario jer menja dan/noc dinamiku.

\subsection{Research gap}
% Treba reci:
% - Postoje fizioloske jednacine i modeli.
% - Problem je sto carbon input cesto ostaje prescribed/simplified.
% - Nedostaje transparentan coupled simulation workflow.

\subsection{Contribution of this work}
% Ovde ide kontribuciona lista.

\begin{enumerate}
    \item Reformulation of a Tedone-based physiological framework as a coupled dynamic simulation model.
    \item Integration of a Farquhar--Jarvis gas-exchange layer.
    \item Propagation of computed assimilation into carbohydrate, nutrient, and biomass dynamics.
    \item Reproducible numerical implementation with defined states, inputs, parameters, and outputs.
    \item Literature-based validation under \textit{Arabidopsis thaliana} photoperiod scenarios.
    \item Ablation and sensitivity analysis of the proposed coupling.
\end{enumerate}

% ============================================================
\section{Related Work}
\label{sec:related_work}
Ovo poglavlje sumira ključne pravce istraživanja o simulacionim modelima biljaka i biofizičkim modulima koji su relevantni za predloženi pristup u ovom radu. Fokus je na mehanističkim modelima rasta, modelima fotosinteze i stomatalne provodljivosti, dinamici ugljenih hidrata, nutritivnim procesima i pristupima koji omogućavaju povezivanje ovih komponenti u jedinstven simulacioni okvir.
% BITNO:
% Kratko. Ne prepisivati dugacki master literature review.
% Digital twin samo eventualno 1 kratak pasus kao future context.

\subsection{Mechanistic plant growth simulation models}
% Sta ide:
% - Mechanistic/process-based models.
% - ODE/PDE/dynamic-system framing.
% - Carbon allocation, nutrient constraints, biomass growth.

Simulacioni modeli biljaka predstavljaju osnovu za povezivanje fizioloških procesa, mikroklimatskih uslova i promena biomase kroz vreme. Prostorne i vremenske skale koje modeli pokrivaju određuju gde i kako se oni mogu primeniti. Prostorno, modeli mogu obuhvatiti nivo lista, organa, cele biljke, sklopa ili čitavog polja. Vremenski, oni pokrivaju strategije od minuta i sati, za mikroklimatske i fiziološke procese, preko dana, za rast i raspodelu ugljenika, do sezone, za procenu prinosa. Pregled funkcionalno-strukturnih modela biljaka ukazuje da modeli koji kombinuju arhitekturu biljke i fiziološke module omogućavaju povezivanje prostorne strukture, fenotipskih informacija i dinamičkih procesa rasta \cite{Louarn2020,Sievnen2014}. Arhitektura savremenih simulacionih modela biljaka obično je modularna. Fotosintetička komponenta može koristiti biohemijski Farquhar-ov model, stomatalna provodljivost može biti opisana preko Ball--Berry modela, Medlyn-ovog optimalističkog pristupa ili Jarvis-ovih funkcija, dok se vodni bilans često povezuje sa Penman--Monteith formulacijom. Takva struktura omogućava pojedinačnu validaciju modula i fleksibilno povezivanje fotosinteze, stomatalne provodljivosti, transpiracije, alokacije ugljenika i rasta \cite{Farquhar1980,Ball1987,Medlyn2011,Jarvis1976,Allen1998,Kim2003}. Modeli kao što su APSIM, DSSAT i WOFOST predstavljaju proverene platforme za simulaciju rasta useva, dinamike vode, azota i klimatskih scenarija \cite{Holzworth2014,Jones2003,deWit2019}. S druge strane, funkcionalno-strukturni modeli i pristupi kao što je GreenLab uvode dodatnu dimenziju povezivanja arhitekture biljke i fizioloških procesa \cite{DeReffye2020}. Ovakvi modeli pokazuju da modularna organizacija ima značajnu prednost kada je potrebno povezati više procesa koji deluju na različitim vremenskim i prostornim skalama. U okviru simulacionih modela biljaka mogu se razlikovati različiti pristupi. Mehanistički, odnosno proces-orijentisani modeli, oslanjaju se na biofizičke i biohemijske jednačine kako bi opisali fotosintezu, transpiraciju, rast i bilanse vode i azota kroz jasno definisane mehanizme. Nedavne studije pokazuju da ovakvi modeli, za razliku od isključivo empirijskih pristupa, bolje odgovaraju varijacijama klimatskih uslova tokom sezone \cite{Mamassi2024}. Noviji primeri uključuju modele kao što je T\&C-CROP, koji integrišu složene poljoprivredne mere, uključujući rotacije useva, navodnjavanje i đubrenje, i pokazuju visok nivo validacije sa realnim podacima \cite{BuckleyPaules2025}. Drugu grupu čine modeli kretanja vode kroz sistem zemljište--biljka. Jedan primer je jednodimenzionalni model koji povezuje ksilem unutar biljke sa zemljištem i pokazuje da ovakav pristup može biti efikasan u poređenju sa složenijim trodimenzionalnim modelima, uz značajno smanjenje računske složenosti \cite{Janott2011}. Posebnu grupu čine diskretni modeli razvoja biljke, kao što su DEM modeli, koji se primenjuju u analizi interakcija korena i zemljišta, uključujući fizičke kontakte između semena, korena i čestica zemljišta \cite{Tian2024}. Modeli optimalnosti zasnovani su na pretpostavci da biljke alociraju resurse na način koji maksimizuje određeni kriterijum, kao što su prinos, rast ili efikasnost resursa. U tom kontekstu, Dynamic Energy Budget teorija predstavlja jedan od okvira za opis raspodele resursa u promenljivim uslovima okruženja \cite{Russo2022}. Konačno, kontinuirani modeli opisuju dinamiku biljke sistemima običnih ili parcijalnih diferencijalnih jednačina. Ovakav pristup omogućava formalnu matematičku analizu sistema i koristi se u modelima u kojima se biomasa, voda ili vegetaciona dinamika predstavljaju kao glatke vremenski promenljive \cite{Meron2019}.


```latex
\subsection{Photosynthesis and stomatal conductance modelling}

Fotosinteza je jedan od centralnih procesa u mehanističkim modelima rasta biljaka, jer određuje primarni ulaz ugljenika koji se kasnije raspodeljuje ka ugljenim hidratima, respiraciji, održavanju i rastu. Za C$_3$ biljke najčešće korišćen biohemijski opis predstavlja Farquhar--von Caemmerer--Berry model, koji povezuje neto asimilaciju sa intercelularnom koncentracijom CO$_2$, Rubisco-limitiranom karboksilacijom i regeneracijom RuBP \cite{Farquhar1980}. Ovaj model se široko koristi kao osnova za opis fotosintetičkog odgovora lista i često se parametrizuje kroz veličine kao što su $V_{cmax}$, $J$, $\Gamma^*$, $K_c$, $K_o$ i $R_d$ \cite{Sharkey2007,Dubois2007}.

Međutim, fotosintetički odgovor ne zavisi samo od biohemijskog kapaciteta lista, već i od difuzione dostupnosti CO$_2$, koja je određena stomatalnom provodljivošću. Stomatalni modeli zato čine neophodnu vezu između mikroklimatskih uslova i asimilacije ugljenika. U literaturi su česti Jarvis-ov empirijski pristup, Ball--Berry model i Medlyn-ov optimalistički model stomatalne regulacije \cite{Jarvis1976,Ball1987,Medlyn2011}. Jarvis-ov pristup opisuje stomatalnu provodljivost kao proizvod funkcija odgovora na spoljašnje faktore, kao što su osvetljenje, deficit pritiska vodene pare, temperatura, koncentracija CO$_2$ i vodni status biljke \cite{Jarvis1976,YU2017,Yu2021}. Takva formulacija je pogodna za simulacione modele jer omogućava direktno uključivanje vremenski promenljivih ulaza iz scenarija ili merenja.

Povezivanje fotosinteze i stomatalne provodljivosti uvodi povratnu spregu između biohemijske potrošnje CO$_2$ i difuzionog toka kroz stome. Intercelularna koncentracija CO$_2$ određuje brzinu asimilacije u Farquhar-ovom modelu, dok stomatalna provodljivost određuje koliko CO$_2$ može da difunduje iz okolnog vazduha u unutrašnjost lista. Zbog toga spoj Farquhar-ovog biohemijskog modela i Jarvis-ove stomatalne regulacije predstavlja prirodan osnov za simulaciju gasne razmene pod promenljivim mikroklimatskim i fotoperiodskim uslovima.

U kontekstu ovog rada, Farquhar-ove i Jarvis-ove jednačine ne predstavljaju samostalan novi fiziološki doprinos, već nasledne modelne komponente koje se integrišu u širi dinamički simulacioni okvir. Doprinos rada nalazi se u njihovom povezivanju sa ugljenohidratnim, nutritivnim i biomasnim promenljivama, tako da se fotosintetski ulaz ne posmatra kao unapred zadat spoljašnji faktor, već kao rezultat simulirane gasne razmene i uslova okruženja.

\subsection{Carbohydrate and photoperiod-dependent growth in \textit{Arabidopsis thaliana}}

Dinamika ugljenih hidrata predstavlja ključnu vezu između fotosinteze i rasta biljke. Tokom svetle faze dana asimilovani ugljenik se raspodeljuje između mobilnih šećera, pre svega saharoze, i privremenih rezervi, pre svega skroba. Tokom tamne faze, kada fotosintetski ulaz prestaje, skrob se razgrađuje i obezbeđuje kontinuirano snabdevanje ugljenikom za respiraciju, održavanje i rast \cite{Stitt2012,Smith2007}. Zbog toga modeli koji opisuju samo trenutnu fotosintezu, bez eksplicitne dinamike skroba i saharoze, imaju ograničenu mogućnost da predstave dan--noć ravnotežu rasta.

Kod \textit{Arabidopsis thaliana}, više radova pokazuje da je noćna degradacija skroba strogo regulisana i da se rezerve skroba često troše približno do trenutka svitanja, čime se smanjuje rizik od noćnog nedostatka ugljenika \cite{Scialdone2013,Stitt2012,Webb2015}. Matematički modeli ovog procesa povezuju cirkadijalnu regulaciju, dostupnost saharoze, aktivnost enzima i signale šećera sa dinamikom rasta i raspodele ugljenika \cite{Feugier2013,Seki2017}. Ova svojstva čine \textit{Arabidopsis thaliana} pogodnim sistemom za validaciju modela koji eksplicitno opisuju fotoperiodski zavisnu dinamiku ugljenih hidrata.

Fotoperiod značajno utiče na odnos između dnevne akumulacije ugljenika i noćne potrošnje rezervi. Kraći dan ograničava vreme fotosintetske asimilacije i produžava period u kome rast zavisi od mobilizacije skroba, dok duži dan omogućava duži period direktnog ugljeničnog ulaza. Eksperimentalni podaci za \textit{Arabidopsis thaliana} pokazuju da biljka koordinira diurnalnu regulaciju alokacije ugljenika i rast kroz širok opseg fotoperioda \cite{Sulpice2014}. Zbog toga fotoperiodski scenariji predstavljaju pogodan validacioni okvir za proveru da li simulacioni model reprodukuje očekivane promene u relativnoj stopi rasta, biomasi i ugljenohidratnim promenljivama.

Pored skroba i saharoze, rast zavisi i od nutritivnog statusa biljke. Azot i fosfor utiču na fotosintetički kapacitet, metabolizam i raspodelu biomase, pa se u mehanističkim modelima često uvode kao unutrašnja stanja ili ograničavajući faktori \cite{Gastal2002,Xu2012,Barber1995,Shen2011,Carstensen2018}. Raspodela ugljenika između lista i korena može se posmatrati kroz odnose izvor--ponor, pri čemu dostupnost asimilata, unutrašnje potrebe za nutrijentima i stanje korena utiču na alokaciju resursa \cite{Marcelis1996,Feller2015,Louarn2020}. U ovom radu se ovi procesi koriste kao deo spregnutog dinamičkog sistema koji povezuje gasnu razmenu, ugljenohidratne rezerve, nutritivne tokove i rast biomase.

\subsection{Simulation validation, ablation, and sensitivity analysis}

U dinamičkim modelima rasta biljaka validacija obično obuhvata poređenje simuliranih i referentnih vrednosti za biomasu, relativnu stopu rasta, tokove vode ili nutrijenata, kao i druge fiziološke promenljive kada su dostupne \cite{Wallach2018,Holzworth2014,Jones2003,deWit2019}. Za model koji se validira pod različitim fotoperiodima, posebno su značajni pokazatelji kao što su relativna stopa rasta, finalna biomasa i dinamika skroba i saharoze.

Pored standardne validacije, analiza ponašanja modela treba da pokaže i doprinos pojedinačnih komponenti. Ablacioni eksperimenti omogućavaju da se uporede varijante modela u kojima se pojedini delovi uklanjaju, pojednostavljuju ili zamenjuju alternativnim formulacijama. U ovom kontekstu, relevantno je porediti osnovni model sa propisanim ili pojednostavljenim fotosintetskim ulazom, model sa Farquhar-ovom fotosintezom, model sa stomatalnom regulacijom i punu varijantu u kojoj su gasna razmena, ugljeni hidrati, nutrijenti i biomasa povezani u jedinstven dinamički sistem. Takvo poređenje pomaže da se pokaže da doprinos rada nije samo implementacija poznatih jednačina, već njihovo funkcionalno povezivanje u simulacionom toku.

Analiza osetljivosti predstavlja dodatni korak u proceni robusnosti modela. Ona pokazuje koji parametri najviše utiču na izlaze modela, kao što su finalna biomasa, relativna stopa rasta, neto asimilacija, skrob i saharoza. U slučaju modela koji uključuje gasnu razmenu, posebno su važni parametri povezani sa svetlošću, temperaturom, ambijentalnom koncentracijom CO$_2$, relativnom vlažnošću ili deficitom pritiska vodene pare, kao i biohemijski parametri $V_{cmax}$ i $J$. Za nizvodne procese važni su parametri degradacije skroba, alokacije saharoze, nutritivnih ograničenja i početnih rezervi ugljenih hidrata.

U ovom radu se validacija, ablation study i sensitivity analysis koriste kao komplementarni elementi evaluacije. Validacija proverava slaganje sa literaturnim ciljevima, ablation study procenjuje doprinos spajanja gasne razmene sa ostatkom modela, dok analiza osetljivosti ispituje stabilnost i dominantne izvore varijabilnosti u ponašanju simulacije.
```


% ============================================================
\section{Model Scope and Contribution Boundary}
\label{sec:scope}

% CILJ:
% Odmah zastititi rad od komentara:
% "ovo je samo implementacija Tedone/Farquhar/Jarvis".
% Ovde ide tabela sta je inherited, sta je adapted, sta je contribution.

\begin{table}[ht]
\centering
\caption{Model components, sources, and contribution boundary.}
\label{tab:component_boundary}
\begin{tabular}{p{0.25\textwidth} p{0.23\textwidth} p{0.32\textwidth} p{0.12\textwidth}}
\hline
\textbf{Component} & \textbf{Source / basis} & \textbf{Role in this work} & \textbf{Status} \\
\hline
Growth and carbohydrate framework & Tedone-based literature &
Baseline carbon storage, allocation, and growth structure & Adopted/adapted \\

Farquhar C$_3$ photosynthesis & Literature &
Assimilation response to $C_i$, light, and biochemical capacity & Integrated \\

Jarvis stomatal conductance & Literature &
Environmental modifiers of stomatal conductance & Integrated \\

Gas-exchange coupling & This work &
Links $g_s$, $C_i$, $A_n$, and downstream carbon input & Proposed coupling \\

Photoperiod simulation workflow & This work &
Controls light/dark forcing and scenario execution & Implemented \\

Validation protocol & This work &
Evaluates model against photoperiod-dependent targets & Designed \\

Ablation and sensitivity analysis & This work &
Tests contribution and robustness of model components & New analysis \\
\hline
\end{tabular}
\end{table}

% ============================================================
\section{Mathematical Model Formulation}
\label{sec:model}

% Ovo je centralna sekcija rada.
% Mora biti formalna, ne opisna kao master.

\subsection{State vector, input vector, and dynamic system}

% LOCKED FORM:
\begin{equation}
    \frac{dX(t)}{dt} = f\left(X(t), U(t), \boldsymbol{\theta}, t\right).
\end{equation}

% LOCKED/LIKELY STATE VECTOR:
\begin{equation}
X(t) =
\left[
k_s,
a,
s,
\alpha_N,
\alpha_P,
N,
P,
\rho_r,
B_l,
B_r
\right]^T .
\end{equation}

% Znacenje:
% k_s      - starch partition coefficient
% a        - starch pool
% s        - sucrose pool
% alpha_N  - nitrogen affinity / investment state
% alpha_P  - phosphorus affinity / investment state
% N        - internal nitrogen state
% P        - internal phosphorus state
% rho_r    - sucrose/root allocation state
% B_l      - leaf biomass
% B_r      - root biomass

% LOCKED/LIKELY INPUT VECTOR:
\begin{equation}
U(t) =
\left[
PAR(t),
T(t),
C_a(t),
RH(t),
VPD(t),
\psi_w(t),
L(t)
\right]^T .
\end{equation}

% Znacenje:
% PAR     - photosynthetically active radiation
% T       - temperature
% C_a     - ambient CO2
% RH/VPD  - humidity / vapour pressure deficit
% psi_w   - water status / water potential
% L(t)    - day/night photoperiod signal

\subsection{Photoperiod forcing}

% LOCKED:
\begin{equation}
L(t) =
\begin{cases}
1, & 0 \leq h(t) < f, \\
0, & f \leq h(t) < 24,
\end{cases}
\end{equation}

% LOCKED if we keep current sinusoidal PAR profile:
\begin{equation}
PAR(t) =
PAR_{\max}
\sin\left(
\pi \frac{h(t)}{f}
\right),
\quad 0 \leq h(t) < f .
\end{equation}

\begin{equation}
PAR(t) = 0, \quad f \leq h(t) < 24 .
\end{equation}

% Sta ovde objasniti:
% - f je photoperiod in hours.
% - h(t) je hour of day.
% - PAR_max je scenario/parameter value.
% - Sinusni profil je idealizovana dnevna kriva.

\subsection{Jarvis stomatal conductance}

% LOCKED GENERAL FORM:
\begin{equation}
g_s =
g_{s,\max}
f_{PAR}(PAR)
f_{VPD}(VPD)
f_T(T)
f_{CO_2}(C_a)
f_w(\psi_w).
\end{equation}

% LOCKED VPD FORM if kept:
\begin{equation}
e_s = 0.61078 \exp\left(\frac{17.27T}{T+237.3}\right),
\end{equation}

\begin{equation}
e_a = e_s \frac{RH}{100},
\end{equation}

\begin{equation}
VPD = e_s - e_a.
\end{equation}

% Sta ovde ide:
% - Definisati sve Jarvis faktore.
% - Ne preterivati sa biologijom.
% - Objasniti da Jarvis daje environment-dependent stomatal conductance.

\subsection{Farquhar C$_3$ photosynthesis}

% LOCKED:
\begin{equation}
A_c =
V_{cmax}
\frac{C_i - \Gamma^*}
{C_i + K_c \left(1 + O/K_o\right)} .
\end{equation}

% LOCKED:
\begin{equation}
A_j =
J
\frac{C_i - \Gamma^*}
{4.5C_i + 10.5\Gamma^*} .
\end{equation}

% LOCKED:
\begin{equation}
A_n = \min(A_c, A_j) - R_d .
\end{equation}

% Sta ovde ide:
% - Ac Rubisco-limited.
% - Aj light-limited / RuBP regeneration-limited.
% - An net assimilation.
% - Svi parametri u posebnoj tabeli.

\subsection{Gas-exchange closure}

% LOCKED GENERAL DIFFUSION RELATION:
\begin{equation}
C_i = C_a - \frac{A_n}{g_c}.
\end{equation}

% LOCKED if gs is water-vapour conductance:
\begin{equation}
g_c \approx \frac{g_s}{1.6}.
\end{equation}

% CONDITIONAL:
% Ako aktiviramo iterate_ci(), ubacujemo ovo kao glavni deo.

\begin{equation}
C_i^{(0)} = \eta C_a, \quad \eta \in [0.6,0.8].
\end{equation}

\begin{equation}
\left| C_i^{(k+1)} - C_i^{(k)} \right| < \varepsilon .
\end{equation}

% Ako NE aktiviramo iterate_ci():
% Ovaj subsection mora da se zove:
% "Jarvis-regulated surrogate approximation of net assimilation"
% i mora da objasni surrogate An = S(PAR, B_l, g_s).
% Tada ne smemo tvrditi full iterative Farquhar--Jarvis runtime.

\subsection{Conversion from assimilation to model carbon input}

% LOCKED GENERAL FORM:
\begin{equation}
PH(t) = \mathcal{C}\left(A_n(t), SLA, B_l(t), \boldsymbol{\theta}\right).
\end{equation}

% Sta ovde mora:
% - Definisati jedinice An.
% - Definisati jedinice PH.
% - Objasniti SLA.
% - Objasniti da PH ulazi u starch/sucrose ODE.
% - Ako ostane base_scale = 1.10, objasniti kao calibration/scaling coefficient.

\subsection{Carbohydrate dynamics}

% LOCKED GENERAL FORM:
\begin{equation}
\frac{da}{dt} =
k_s PH - D_a .
\end{equation}

\begin{equation}
\frac{ds}{dt} =
(1-k_s)PH + D_a
- C_{uptake}
- C_{transport}
- C_{resp}
- C_{load}.
\end{equation}

% Sta ovde ide:
% - a = starch.
% - s = sucrose.
% - D_a = starch degradation.
% - day: PH > 0.
% - night: PH = 0, starch degradation supplies sucrose.

% LOCKED/LIKELY:
\begin{equation}
PH(t) = 0, \quad L(t)=0 .
\end{equation}

\subsection{Nutrient uptake and nutrient-related costs}

% LOCKED/LIKELY:
\begin{equation}
U_N^{pot} =
I_{N,\max}
\frac{N_s}{N_s + K_N}
\alpha_N
\frac{B_r}{B_l+B_r}.
\end{equation}

\begin{equation}
U_P^{pot} =
I_{P,\max}
\frac{P_s}{P_s + K_P}
\alpha_P
\frac{B_r}{B_l+B_r}.
\end{equation}

\begin{equation}
U_N =
U_N^{pot}
\frac{s}
{s + c_N U_N^{pot} + \epsilon}.
\end{equation}

\begin{equation}
U_P =
U_P^{pot}
\frac{s}
{s + c_P U_P^{pot} + \epsilon}.
\end{equation}

% Sta ovde ide:
% - Potential uptake depends on soil availability, affinity, root biomass share.
% - Actual uptake is sucrose-limited.
% - N/P dynamics consume carbon and affect growth/photosynthetic capacity.

\subsection{Biomass growth and allocation}

% LOCKED GENERAL FORM:
\begin{equation}
\frac{dB_l}{dt} =
G_l(s, \rho_r, \boldsymbol{\theta}) - M_l(B_l).
\end{equation}

\begin{equation}
\frac{dB_r}{dt} =
G_r(s, \rho_r, \boldsymbol{\theta}) - M_r(B_r).
\end{equation}

% Sta ovde ide:
% - Leaf/root growth depends on sucrose availability.
% - rho_r controls allocation to root.
% - Mortality/competition terms included.
% - Full exact implementation can be listed after final code check.

% ============================================================
\section{Algorithmic Description}
\label{sec:algorithm}

% Ovde ide Algorithm 1.
% Ne kao software detail, nego kao modelling/simulation contribution.

\begin{algorithm}
\caption{One simulation step with Farquhar--Jarvis gas-exchange coupling}
\label{alg:simulation_step}
\begin{algorithmic}[1]
\STATE Read $X(t)$, $U(t)$, $\boldsymbol{\theta}$, and $\Delta t$.
\STATE Determine $L(t)$ from the photoperiod function.
\STATE Compute or read $PAR(t)$, $T(t)$, $C_a(t)$, $RH(t)$ or $VPD(t)$, and water status.
\STATE Compute Jarvis environmental modifiers.
\STATE Compute stomatal conductance $g_s$.
\STATE Compute or solve $C_i$.
\STATE Compute $A_c$, $A_j$, and $A_n$.
\STATE Convert $A_n$ into model carbon input $PH$.
\STATE Update starch, sucrose, nutrient-related states, allocation state, and biomass states.
\STATE Store outputs and diagnostics.
\end{algorithmic}
\end{algorithm}

% CONDITIONAL:
% If iterate_ci() is active, expand steps 6-7 into iterative loop.
% If surrogate remains active, explicitly say "compute An using surrogate approximation".

% ============================================================
\section{Numerical Implementation}
\label{sec:implementation}

% Sta ovde ide:
% - C numerical core.
% - External solver framework.
% - RHS function.
% - Algebraic variables computed before derivatives.
% - CSV input scenario.
% - Output time series.
% - Validation runner.
% - Reproducibility: config, commit hash, solver version.

\subsection{Implementation architecture}
% Mention modules:
% - photosynthesis
% - carbohydrates
% - nutrients
% - growth
% - algebraic outputs
% - simulation dynamics
% - validation runner

\subsection{Numerical integration}
% TODO:
% - exact solver used
% - dt_hours
% - tolerances
% - fixed/adaptive step
% - time-step sensitivity check

\subsection{Input and output workflow}
% Sta ide:
% - Input CSV scenario.
% - Environmental time series.
% - Initial conditions.
% - Parameter files.
% - Output CSV with states and diagnostics.

\subsection{Reproducibility}
% Sta ide:
% - Build command.
% - Commit hash.
% - solver.lock / solver version.
% - validation runner command.
% - output CSV files.

% ============================================================
\section{Simulation Scenarios and Validation Protocol}
\label{sec:validation}

% Sta ovde ide:
% - Validation is literature-based.
% - Arabidopsis is validation domain.
% - Photoperiods: 4h, 6h, 8h, 12h.
% - Metrics: RGR error, MAE, RMSE, nRMSE, pass/fail.
% - State validity checks.
% - Day/night photosynthesis checks.
% - Time-step sensitivity.

\subsection{Photoperiod scenarios}

\begin{table}[ht]
\centering
\caption{Photoperiod validation scenarios.}
\label{tab:scenarios}
\begin{tabular}{rrrr}
\hline
\textbf{Scenario} & \textbf{Photoperiod} & \textbf{Duration} & \textbf{Primary target} \\
\hline
S1 & 4 h  & TODO & RGR \\
S2 & 6 h  & TODO & RGR, sugars/starch if available \\
S3 & 8 h  & TODO & RGR \\
S4 & 12 h & TODO & RGR \\
\hline
\end{tabular}
\end{table}

\subsection{RGR validation}

% LOCKED:
\begin{equation}
RGR =
\frac{\ln B(t_2) - \ln B(t_1)}
{t_2 - t_1}.
\end{equation}

% Or if implementation uses start/final ratio differently, align with code.

\begin{table}[ht]
\centering
\caption{RGR validation under photoperiod treatments.}
\label{tab:rgr_validation}
\begin{tabular}{rrrrr p{0.22\textwidth}}
\hline
\textbf{Photoperiod} & \textbf{Literature RGR} & \textbf{Simulated RGR} &
\textbf{Abs. error} & \textbf{Rel. error} & \textbf{Comment} \\
\hline
4 h  & 0.0680 & TODO & TODO & TODO & TODO \\
6 h  & 0.1135 & TODO & TODO & TODO & TODO \\
8 h  & 0.1708 & TODO & TODO & TODO & TODO \\
12 h & 0.2600 & TODO & TODO & TODO & TODO \\
\hline
\end{tabular}
\end{table}

\subsection{Sugar and starch validation}
In addition to RGR, the model output includes simulated sucrose and starch
pools, which are used to evaluate whether the model reproduces physiologically
plausible day–night carbohydrate dynamics. When a curated reference dataset is
available, sucrose and starch are compared quantitatively against literature-derived
 time points. The comparison is performed separately for each photoperiod
scenario for which reference values are available.
The following metrics are used:
\begin{equation}
MAE =
\frac{1}{n}
\sum_{i=1}^{n}
\left|
y_i - \hat{y}_i
\right|,
\end{equation}
\begin{equation}
RMSE =
\sqrt{
\frac{1}{n}
\sum_{i=1}^{n}
\left(
y_i - \hat{y}_i
\right)^2
},
\end{equation}
\begin{equation}
nRMSE =
\frac{RMSE}
{y_{\max} - y_{\min}}.
\end{equation}
Here, $y_i$ denotes the reference value, $\hat{y}
_i$ denotes the simulated value at the corresponding time point, and $n$ is the
number of available reference points. When reference values are available only
in relative or normalized form, the comparison is performed using normalized
time-series behaviour rather than absolute concentrations.

\begin{table}[ht]
\centering
\caption{Sugar and starch validation summary.}

\label{tab:sugar_starch_validation}
\begin{tabular}{rrrrr p{0.22\textwidth}}
\hline
\textbf{Photoperiod} &
\textbf{Ref. points} &
\textbf{Sucrose nRMSE} &
\textbf{Sucrose MAE} &
\textbf{Starch nRMSE} &
\textbf{Starch MAE} &
\textbf{Comment} \\
\hline
4 h & TODO & TODO & TODO & TODO & TODO & TODO \\
6 h & TODO & TODO & TODO & TODO & TODO & TODO \\
8 h & TODO & TODO & TODO & TODO & TODO & TODO \\
12 h & TODO & TODO & TODO & TODO & TODO & TODO \\
\hline
\end{tabular}
\end{table}

If a sufficient number of curated reference points is not available for a given
photoperiod, sugar and starch are not treated as fully quantitative validation
targets for that scenario. In such cases, the model is evaluated qualitatively
with respect to expected carbohydrate behaviour: starch should accumulate during
the light phase, decrease during the dark phase, and contribute to sucrose
availability during the night. This distinction is reported explicitly to avoid
interpreting trend-level agreement as full quantitative validation

\subsection{Gas-exchange validation or plausibility check}
% Sta ide:
% - If direct An/gs/Ci data exist: quantitative validation.
% - If not: range and trend plausibility.
% - If surrogate remains: surrogate error must be validated.

\subsection{Numerical checks}
% Sta ide:
% - finite states
% - non-negative sucrose/starch/biomass
% - positive final biomass
% - night photosynthesis close to zero
% - dt vs dt/2 sensitivity

% ============================================================
\section{Ablation Study}
\label{sec:ablation}

% Cilj:
% Pokazati sta konkretno menja predlozeni coupling.

\begin{table}[ht]
\centering
\caption{Ablation variants.}
\label{tab:ablation_variants}
\begin{tabular}{rrrrr p{0.22\textwidth}}
\hline
\textbf{Variant} & \textbf{Description} \\
\hline
M0 & Tedone-style baseline with simplified or prescribed carbon input. \\
M1 & Tedone-based model with Farquhar photosynthesis only. \\
M2 & Tedone-based model with Farquhar--Jarvis gas-exchange coupling. \\
M3 & Full model with photoperiod forcing, gas exchange, carbohydrates, nutrients, and biomass dynamics. \\
\hline
\end{tabular}
\end{table}

% Metrics:
% - RGR error
% - final biomass error
% - sucrose/starch dynamics
% - day/night plausibility
% - numerical stability

% ============================================================
\section{Sensitivity Analysis}
\label{sec:sensitivity}

% Minimum: one-at-a-time.
% Better: Morris/Sobol if time allows.

\subsection{Parameters}
% Candidate parameters:
% - PAR_max
% - T
% - C_a
% - RH/VPD
% - Vcmax
% - J
% - Jarvis VPD sensitivity
% - starch degradation rate
% - sucrose allocation coefficient
% - nutrient limitation coefficient
% - initial starch/sucrose

\subsection{Outputs}
% Candidate outputs:
% - final biomass
% - RGR
% - An / PH
% - sucrose
% - starch
% - N
% - P

\subsection{Sensitivity metric}
% TODO:
% Choose metric.
% Simple option:
% relative output change per relative parameter perturbation.

\begin{equation}
S_p =
\frac{\Delta Y / Y_0}
{\Delta p / p_0}.
\end{equation}

% ============================================================
\section{Results}
\label{sec:results}

% Ne pisati jos dok nema finalnih runova.

\subsection{Numerical behaviour of the gas-exchange coupling}
% Sta ide:
% - convergence if iterate_ci active
% - or surrogate behaviour if surrogate path active
% - diagnostics

\subsection{Photoperiod-dependent assimilation response}
% Sta ide:
% - PAR/An/PH through day
% - comparison across 4h, 6h, 8h, 12h

\subsection{Carbohydrate dynamics}
% Sta ide:
% - starch accumulation in day
% - starch degradation at night
% - sucrose behaviour

\subsection{Biomass and RGR validation}
% Sta ide:
% - RGR table
% - final biomass
% - errors

\subsection{Ablation results}
% Sta ide:
% - M0-M3 comparison

\subsection{Sensitivity results}
% Sta ide:
% - tornado plot / table
% - dominant parameters
% - unstable regions

% ============================================================
\section{Discussion}
\label{sec:discussion}

% Sta mora eksplicitno:
% - This is not a completely new plant physiology model.
% - Tedone/Farquhar/Jarvis are inherited and cited.
% - Contribution is coupling, reformulation, implementation, validation.
% - Arabidopsis is validation domain.
% - This is not a full digital twin.
% - No new wet-lab validation is claimed.
% - Sensor integration, actuator feedback, online correction, surrogate acceleration are future work.

% ============================================================
\section{Limitations}
\label{sec:limitations}

% Moze skoro fiksno:
% - literature-parameterized model
% - limited published validation data
% - simplified root architecture
% - simplified hormonal/developmental regulation
% - no full leaf energy balance
% - parameter uncertainty
% - no complete digital twin deployment
% - surrogate not central unless validated
% - transfer to other species requires recalibration

% ============================================================
\section{Conclusion}
\label{sec:conclusion}

% Sta mora:
% - Presented extended Tedone-based mechanistic simulation model.
% - Integrated Farquhar--Jarvis gas-exchange layer.
% - Linked environmental drivers to carbon assimilation and downstream growth.
% - Validated under Arabidopsis photoperiod treatments.
% - Ablation and sensitivity analysis support relevance/robustness.
% - Provides simulation core for future hybrid/digital-twin-enabled work.
```



% Numbered list
% Use the style of numbering in square brackets.
% If nothing is used, default style will be taken.
%\begin{enumerate}[a)]
%\item 
%\item 
%\item 
%\end{enumerate}  

% Unnumbered list
%\begin{itemize}
%\item 
%\item 
%\item 
%\end{itemize}  

% Description list
%\begin{description}
%\item[]
%\item[] 
%\item[] 
%\end{description}  

\clearpage %%Remove this from your manuscript

% Figure
\begin{figure}%[]
  \centering
%    \includegraphics{}
    \caption{}\label{fig1}
\end{figure}


\begin{table}%[]
\caption{}\label{tbl1}
\begin{tabular*}{\tblwidth}{@{}LL@{}}
\toprule
  &  \\ % Table header row
\midrule
 & \\
 & \\
 & \\
 & \\
\bottomrule
\end{tabular*}
\end{table}

% Uncomment and use as the case may be
%\begin{theorem} 
%\end{theorem}

% Uncomment and use as the case may be
%\begin{lemma} 
%\end{lemma}

%% The Appendices part is started with the command \appendix;
%% appendix sections are then done as normal sections
%% \appendix

\section{}\label{}

% To print the credit authorship contribution details
\printcredits

%% Loading bibliography style file
%\bibliographystyle{model1-num-names}
\bibliographystyle{cas-model2-names}

% Loading bibliography database
\bibliography{master_ref}

% Biography
%\bio{}
% Here goes the biography details.
%\endbio

%\bio{pic1}
% Here goes the biography details.
%\endbio

\end{document}

